% Modernised reading — fonds Alexandre Grothendieck, shelfmark 113, pages 1-12.
%
% This is an INTERPRETATION, not a transcription. Everything the critical
% apparatus carried moves into footnotes. In any dispute of substance the
% transcription governs: whoever cites Grothendieck must cite batch-01.fr.tex.
%
% Pass: Opus 5 (claude-opus-5), 2026-09-13 — first pass, unchecked against the
% pages by a human. The folder was read whole; it holds one batch, pages 1-12.
% Page 1 is a cover (« Abelianization »), pages 4, 6, 8, 10 listings with
% nothing of his. Page 2 treats another question and is read briefly.
%
% What the folder is. Under his heading « (x) dans Cat », pages 3-12 build
% the tensor product of small k-linear categories and show that it computes
% the tensor product of their categories of modules: P^k = Hom_k(P°, Ab_k),
% the recognition of such categories, the free-cocompletion adjunction, the
% 2-universal P (x)_k Q with its explicit construction, (P (x)_k Q)^k as the
% tensor product of P^k and Q^k among cocomplete k-linear categories, the
% comparison functor P^k (x)_k M -> P^k_M, projectives in P^k_M, and the
% derived convolution.
%
% Notation fixed for the whole document. k a commutative ring, Ab_k the
% k-modules. For P small k-linear: P^k = Hom_k(P°, Ab_k); for M k-linear,
% P^k_M = Hom_k(P°, M) (his script P_M). Hom_{k!} = k-linear functors
% preserving small colimits; Hom^!_k = those sending colimits to limits.
% Bil_k bilinear functors, Bil_{k!!} those cocontinuous in each variable.
% F *_k L' the convolution (coend) of F in P^k_M with L' in (P°)^k; Hom_k(L, F)
% the corresponding end, in M. « Stable par lim » = admitting small colimits.
%
% Substantive departures from the pages, each footnoted where it occurs:
% — page 3: a) and b) do not give a generating family of SMALL projectives,
%   which is what the reconstruction uses (Freyd-Mitchell); the reading
%   states the criterion with it, and reads « Ul Proj » as that subcategory.
% — page 5: both sides of the first display are Q^k_M = Hom_k(P, M); the
%   leaf's labels script-Q_M and script-P_M do not match the convention of
%   page 9, and the reading uses page 9's.
% — pages 9 and 11: (⊛) P^k (x)_k M -> P^k_M, and its case
%   P^k (x)_k Q^k -> (P (x)_k Q)^k, are NOT fully faithful in general
%   (P = k a field, M = Mod_k, infinite-dimensional objects); they are on
%   L, L' in Kar P when the Hom-modules of P are finitely generated projective.
%   The page asserts it, then ends on « pl. fid. ? ».
% — page 11: the middle term Hom_M(x, Hom_k(L, F)) needs limits in M; the
%   outer isomorphism holds without.
% — page 12: « M |-> M (x)_k a » read with x.
%
% Modern names given and footnoted as ours: Karoubi envelope and Morita
% invariance; Freyd-Mitchell recognition of functor categories; free
% cocompletion; tensor product of k-linear categories; tensor product of
% cocomplete k-linear categories (Kelly), and Deligne's tensor product of
% abelian categories; coend and end; Tor over a ring with several objects.
%
% What the folder announces and never establishes: the 2-functorial
% statement of page 3 (« 2-ess. surj. », « 3-fid. »), only drawn; the
% question of page 11 in its Tor part; the propositions of page 12 beyond the
% two-line arguments given here.
\documentclass[11pt,a4paper]{article}
% Le préambule commun à toutes les transcriptions du fonds Grothendieck.
%
% Il tient en une idée : **l'incertitude se compose, elle ne se résout pas.**
% Une transcription qui lisse un mot douteux a détruit la seule chose pour
% laquelle on la faisait — un lecteur doit pouvoir savoir, à chaque endroit,
% ce qui était sur le feuillet et ce qui a été ajouté. Les six macros qui
% suivent sont donc l'appareil critique, et non de la décoration.
%
% Les mêmes commandes sont reconnues par `scripts/render.mjs`, qui produit la
% vue de lecture du site. Une macro ajoutée ici doit l'être aussi là-bas,
% faute de quoi le rendu échouera bruyamment — ce qui est voulu : un rendu
% silencieusement faux est pire qu'un rendu absent.

\ProvidesPackage{grothendieck}[2026/08/08 Transcriptions du fonds Grothendieck]

\usepackage{amsmath,amssymb,amsthm}
\usepackage{mathrsfs}
\usepackage{stmaryrd}
% Les diagrammes commutatifs sont partout dans ce fonds : c'est la langue
% même de la géométrie algébrique de Grothendieck, et une transcription qui
% les rendrait en prose ne transcrirait rien.
\usepackage{tikz-cd}
% Français par défaut — c'est la langue des feuillets — mais l'anglais est
% chargé aussi : la traduction et le résumé sont des documents anglais, et la
% typographie française (espaces avant les deux-points) y serait une faute.
% Ces documents basculent par \selectlanguage{english} en tête de corps.
\usepackage[english,french]{babel}
\usepackage{fontspec}
\usepackage{geometry}
\geometry{a4paper,margin=25mm,marginparwidth=28mm}
\usepackage{marginnote}
\usepackage{xcolor}
% ulem, not soul. `soul`'s \st and \ul are built on a character-by-character
% scan that XeLaTeX's font machinery breaks: the document fails with an
% undefined control sequence rather than degrading. ulem does the same two jobs
% and is Unicode-engine safe. `normalem` stops it hijacking \emph, which the
% transcriptions use for Grothendieck's own emphasis.
\usepackage[normalem]{ulem}
\usepackage{hyperref}
\hypersetup{colorlinks=true,linkcolor=black,urlcolor=black,pdfborder={0 0 0}}

% --- Métadonnées du lot -----------------------------------------------------
% Reprises telles quelles par le script de rendu, qui les affiche en tête de
% la vue de lecture. Elles fixent ce que le fichier prétend couvrir : sans
% elles, une transcription flotte sans point d'attache dans un fonds de seize
% mille feuillets.

\newcommand{\folder}[1]{\gdef\grothFolder{#1}}
\newcommand{\batch}[1]{\gdef\grothBatch{#1}}
\newcommand{\leaves}[2]{\gdef\grothFirst{#1}\gdef\grothLast{#2}}
\newcommand{\foldertitle}[1]{\gdef\grothFoldertitle{#1}}
\gdef\grothFolder{?}\gdef\grothBatch{?}\gdef\grothFirst{?}\gdef\grothLast{?}\gdef\grothFoldertitle{}

% --- Le résumé --------------------------------------------------------------
%
% Il ouvre la lecture modernisée, sous le titre « Résumé », et il est composé
% en petit. Ce n'est pas un ornement : c'est le seul passage du document qui
% ne soit pas des mathématiques, et un lecteur doit pouvoir le reconnaître
% avant de l'avoir lu — puis le sauter s'il connaît déjà le sujet, ou s'y
% arrêter s'il ne le connaît pas. Le filet qui le suit marque la reprise du
% corps.

\newenvironment{resume}
  {\small\setlength{\parskip}{0.45em}}
  {\par\medskip\hrule\medskip}

% --- Mots-clés ---------------------------------------------------------------
%
% En anglais, en clôture du résumé de la lecture modernisée : le vocabulaire
% moderne sous lequel chercher ce que les feuillets construisent. Le manifeste
% du site (`npm run manifest`) les extrait pour étiqueter la cote entière —
% cette ligne est la source unique des tags, il n'y a pas de fichier à côté.
\newcommand{\keywords}[1]{\par\smallskip\noindent
  {\footnotesize\textsc{Keywords} --- \textit{#1}}\par}

% --- L'appareil critique ----------------------------------------------------

% Le numéro de feuillet, en marge. C'est l'ancre qui permet de régler un
% désaccord sur une formule en regardant *un* feuillet plutôt qu'une chemise
% de six cents. Le site s'en sert aussi pour tourner les pages du fac-similé
% quand on fait défiler la transcription.
\newcommand{\leaf}[1]{%
  \marginnote{\footnotesize\color{gray!70}\textsf{#1}}%
  \label{leaf:#1}\ignorespaces}

% Illisible : on ne devine pas. Un mot inventé qui se lit comme les autres est
% le pire résultat possible.
\newcommand{\ill}{\textcolor{red!60!black}{[\,\ldots\,]}}

% Lecture douteuse : le mot est donné, et signalé comme incertain.
\newcommand{\uncertain}[1]{\uline{#1}}

% Ajout de l'éditeur — une lettre manquante, un mot sous-entendu.
\newcommand{\add}[1]{\textcolor{blue!50!black}{[#1]}}

% Biffé par Grothendieck lui-même. Ce qu'il a rayé fait partie du document :
% une rature dit souvent où la pensée a changé de direction.
\newcommand{\struck}[1]{\sout{#1}}

% Note du transcripteur, sur l'état matériel du feuillet.
\newcommand{\note}[1]{\par\smallskip{\small\color{gray!60!black}\textit{[#1]}}\par\smallskip}

% Note marginale de Grothendieck, distincte du corps du texte.
\newcommand{\marginal}[1]{\par\smallskip\noindent\hspace{1em}%
  {\small\color{gray!60!black}#1}\par\smallskip}

% --- Titre ------------------------------------------------------------------

\AtBeginDocument{%
  \begin{center}
    {\large\bfseries Fonds Alexandre Grothendieck --- cote n\textsuperscript{o}~\grothFolder}\\[2pt]
    {\itshape\grothFoldertitle}\\[4pt]
    {\small feuillets \grothFirst--\grothLast\ (lot \grothBatch)}\\[2pt]
    {\footnotesize\color{gray!60!black}Transcription établie d'après le fac-similé numérisé
     par l'Université de Montpellier.\\
     Les lectures incertaines sont soulignées, les illisibles notées [\,\ldots\,],
     les ajouts entre crochets.}
  \end{center}
  \vspace{1em}}

\endinput


\folder{113}
\pages{1}{12}
\dating{[à partir de 1982]}
\watermark{Édition de démonstration\\interprétation personnelle de l'œuvre}
\foldertitle{Abelianization : notes manuscrites (s.d.) — lecture modernisée du dossier entier}

\begin{document}

\section*{Résumé}

\begin{resume}
Un anneau se connaît par ses modules. On peut aller plus loin et remplacer
l'anneau par une petite catégorie « linéaire » --- une catégorie dont les
ensembles de flèches sont des modules et où la composition est bilinéaire, un
anneau à plusieurs objets en somme --- et considérer ses modules, qui sont les
foncteurs à valeurs dans les modules. Ces catégories de modules sont de grosses
catégories abéliennes, avec toutes les sommes, et l'on sait les reconnaître
parmi toutes les catégories abéliennes.

Pour les anneaux commutatifs, il existe une opération familière : le produit
tensoriel $A \otimes_{k} B$, et les modules sur $A \otimes_{k} B$ sont les objets
qui sont à la fois des $A$-modules et des $B$-modules de manière compatible. Les
feuillets, sous le titre de sa main « $\otimes$ dans Cat », font la même chose
pour les catégories linéaires : ils fabriquent un produit tensoriel
$P \otimes_{k} Q$, par une propriété universelle puis par une construction
explicite, et montrent que la catégorie de modules de $P \otimes_{k} Q$ est le
produit tensoriel des catégories de modules de $P$ et de $Q$ --- au sens
précis où un foncteur bilinéaire sur les deux, qui respecte les sommes dans
chaque variable, est la même chose qu'un foncteur sur le produit.

La preuve est une suite de « curryfications » : un foncteur de deux variables
est un foncteur d'une variable à valeurs dans les foncteurs de l'autre. C'est
d'une grande économie, et c'est la manière moderne de traiter ces questions.

Les dernières pages essaient d'aller plus loin --- comparer les deux produits
tensoriels possibles, construire des résolutions projectives, dériver le
produit --- et c'est là que le texte s'enhardit un peu trop : une comparaison
qu'il déclare pleinement fidèle ne l'est pas en général, parce que le produit
tensoriel de deux espaces d'applications linéaires est plus petit que l'espace
des applications du produit dès que la dimension est infinie. Le feuillet
lui-même finit cette ligne par un point d'interrogation. Au dos du feuillet de
titre, sans rapport, un calcul sur des « MM-fibrés » semi-simpliciaux.

\keywords{k-linear category, functor category, Karoubi envelope, Morita invariance, free cocompletion, tensor product of linear categories, Kelly tensor product, coend, derived tensor product}
\end{resume}

\section*{Le fil du dossier, et les conventions}

\begin{enumerate}
\item[1.] \emph{Au dos du titre} (page 2) : un calcul indépendant.
\item[2.] \emph{Catégories de modules} (page 3) : $P^{k}$, sa reconnaissance,
l'adjonction de la cocomplétion libre.
\item[3.] \emph{Produit tensoriel de catégories $k$-linéaires} (page 5), et
son plongement dans les foncteurs bilinéaires (page 7).
\item[4.] \emph{Trois applications} (pages 7 et 9) : $P^{k}_{M}$ comme catégorie
de modules, $(P \otimes_{k} Q)^{k}$ comme produit tensoriel, et la comparaison
$(\circledast)$.
\item[5.] \emph{Projectifs et produit dérivé} (pages 11 et 12).
\end{enumerate}

\emph{Conventions.} $k$ est un anneau commutatif et $Ab_{k}$ la catégorie des
$k$-modules. Une catégorie $k$-additive est une catégorie additive $k$-linéaire.
Pour $P$ petite et $k$-additive,
\[
P^{k} = \mathrm{Hom}_{k}(P^{\circ}, Ab_{k}),
\qquad
P^{k}_{M} = \mathrm{Hom}_{k}(P^{\circ}, M)
\]
pour $M$ $k$-additive.\footnote{Il écrit $\mathcal{P}$, d'un $P$ bouclé, pour
la grosse catégorie et $P$ pour la petite, et $\mathcal{P}_{M}$ pour
$P^{k}_{M}$ ; on écrit $P^{k}$ et $P^{k}_{M}$. Une catégorie « stable par
$\varinjlim$ », abrégé « st.bl. », est une catégorie qui admet les petites
limites inductives.} $\mathrm{Hom}_{k!}$ désigne les foncteurs $k$-linéaires
qui commutent aux petites limites inductives, $\mathrm{Hom}^{!}_{k}$ ceux qui les
transforment en limites projectives, $\mathrm{Bil}_{k}$ les foncteurs
$k$-bilinéaires et $\mathrm{Bil}_{k!!}$ ceux qui commutent aux limites
inductives en chaque variable. Pour $F \in P^{k}_{M}$ et $L' \in
(P^{\circ})^{k}$, $F *_{k} L' \in M$ est leur convolution, et pour $L \in P^{k}$,
$\mathrm{Hom}_{k}(L, F) \in M$ l'objet des morphismes :
\[
F *_{k} L' = \int^{a \in P} L'(a) \otimes_{k} F(a),
\qquad
\mathrm{Hom}_{k}(L, F) = \int_{a \in P} \mathrm{Hom}_{k}\bigl(L(a), F(a)\bigr),
\]
la première supposant $M$ cocomplète, la seconde complète.\footnote{Les
formules par cofin et fin sont les nôtres ; la page nomme les deux opérations
sans les définir. L'objet $\mathrm{Hom}_{k}(L(a), F(a))$ de $M$ est la
puissance de $F(a)$ par le $k$-module $L(a)$.}

\emph{Ce que le dossier annonce et n'établit pas} : l'énoncé 2-fonctoriel de la
page 3, seulement dessiné ; la partie « Tor » de la question de la page 11 ; les
propositions de la page 12 au-delà des deux arguments donnés ici.

\pagerange{2}{2}
\section*{Au dos du feuillet de titre (page 2)}

Le feuillet de couverture ne porte que le mot « Abelianization » et le nombre
377 ; c'est de lui que vient le titre de l'inventaire. Son verso, écrit en
travers, traite d'une autre question. Des diagrammes
$\alpha \twoheadrightarrow \beta \hookleftarrow \gamma$ au-dessus de
$X \twoheadrightarrow Y \hookleftarrow Z$ en composent une chaîne, et une
famille d'objets notés $M(\xi_{0}, \ldots, \xi_{n})$, reliés par des flèches
presque toutes marquées « subm », décrit ce qu'il appelle des
\emph{MM-fibrés}. Le feuillet se termine sur un critère :

\emph{Pour que $\int_{I} X(i)$ ne fasse pas sortir des MM-fibrés
semi-simpliciaux, il suffit que}
\begin{enumerate}
\item[a)] le nerf de $I$ soit asphérique, c'est-à-dire que $I$ soit asphérique ;
\item[b)] chaque $X(i)$ transforme les opérateurs de face $d_{i} : \Delta_{n-1}
\to \Delta_{n}$ en submersions ;
\item[c)] pour toute flèche $u : i \to j$ de $I$, $X(u)$ se factorise en
$g f$, avec $f$ une submersion et $g$ un monomorphisme.
\end{enumerate}

Cette lecture n'identifie pas la notion.\footnote{« MM-fibré »
remplace « multibundle », biffé. Le mot qualifiant $f$ en c) est lu
« submersion » sur la foi de l'étiquette « subm » qui court sur tout le
feuillet. Une première version du bloc central, biffée en noir puis en rouge,
n'est pas reprise, et la transcription n'a pas pu séparer les extrémités des
flèches obliques $\eta_{0}, \ldots, \eta_{n-1}$. Le listing est daté du 16
septembre 1982.}

\pagerange{3}{3}
\section*{Catégories de modules (page 3)}

Soit $P$ une petite catégorie $k$-additive. La catégorie $P^{k}$ ne dépend pas
de la structure $k$-linéaire choisie : un foncteur additif $P^{\circ} \to Ab$ est
automatiquement à valeurs dans les $k$-modules. Le plongement de Yoneda
$P \hookrightarrow P^{k}$ est additif et pleinement fidèle, et il induit
\[
\mathrm{Kar}(P) \xrightarrow{\;\sim\;} \mathrm{Proj}_{\mathrm{tf}}(P^{k}),
\qquad
(\mathrm{Kar}\, P)^{k} \xrightarrow{\;\sim\;} P^{k},
\]
où $\mathrm{Proj}_{\mathrm{tf}}(P^{k})$ est la sous-catégorie pleine des objets
projectifs de type fini : les facteurs directs des objets représentables.
C'est l'invariance de Morita.\footnote{La page écrit « $\mathrm{Kar}\, P
\approx \mathrm{Ul}\,\mathrm{Proj}(P^{k})$ », les deux lettres devant
« Proj » étant lues « Ul » et non identifiées. Les projectifs de $P^{k}$ ne sont
pas tous dans l'image de $\mathrm{Kar}\, P$ --- une somme infinie de
représentables ne l'est pas --- et c'est la sous-catégorie des projectifs de
type fini que l'énoncé demande ; on lit ainsi. La seconde équivalence est en
marge.}

\textbf{Reconnaissance.} Pour qu'une catégorie $k$-additive $\mathcal{P}$ soit
équivalente à $P^{k}$ pour une petite catégorie $k$-additive $P$, il faut et il
suffit qu'elle soit abélienne, qu'elle admette les petites sommes, et qu'elle
admette une petite famille génératrice d'objets projectifs \emph{de type fini},
c'est-à-dire dont le foncteur $\mathrm{Hom}$ commute aux sommes. On peut alors
prendre pour $P$ la sous-catégorie pleine de ces projectifs.\footnote{La page
demande a) que $\mathcal{P}$ soit un « topos abélien » --- abélienne, stable
par petites limites inductives filtrantes, celles-ci étant exactes, avec une
petite sous-catégorie génératrice --- et b) qu'elle ait assez de projectifs.
Ces conditions donnent bien une famille génératrice de projectifs, mais rien
ne dit que ceux-ci soient de type fini, et c'est ce que la reconstruction
utilise : la page prend elle-même « $P = \mathrm{Ul}\,\mathrm{Proj}(\mathcal{P})$ ».
Le critère juste est dû à P.~Freyd et à B.~Mitchell ; le nom n'est pas sur
la page.}

\emph{Morphismes.} Les morphismes de « topos $k$-abéliens algébriques » sont
les foncteurs $f_{!} : \mathcal{P} \to \mathcal{Q}$ qui commutent aux limites
inductives et dont l'adjoint à droite $f^{*}$ y commute aussi, c'est-à-dire
admet lui-même un adjoint à droite $f_{*}$. Il en dessine le résultat : le
2-foncteur $P \mapsto P^{k}$, de la 2-catégorie des petites catégories
$k$-additives vers celle des topos $k$-abéliens algébriques, est
2-essentiellement surjectif, et devient une équivalence sur les catégories de
morphismes après passage à l'enveloppe de Karoubi.\footnote{Le diagramme porte
« 2-ess.\ surj. » et « 3-fid. » ; on lit le second comme la pleine fidélité au
niveau des 2-flèches. Un tel $f_{!}$ préserve les projectifs de type fini,
puisque $f^{*}$ est exact et commute aux sommes, et c'est ce qui le ramène à
un foncteur $\mathrm{Kar}\,P \to \mathrm{Kar}\,Q$. La page ne démontre rien.}

\emph{Cocomplétion libre.} Tout $L \in P^{k}$ est limite inductive canonique de
représentables, $L \simeq \varinjlim_{P/L} a$, d'où, pour $M$ $k$-additive
admettant les petites limites inductives,
\[
\mathrm{Hom}_{k!}(P^{k}, M) \xrightarrow{\;\sim\;} \mathrm{Hom}_{k}(P, M),
\]
et, dualement, pour $M$ admettant les petites limites projectives,
\[
\mathrm{Hom}^{!}_{k}\bigl((P^{k})^{\circ}, M\bigr) \xrightarrow{\;\sim\;}
\mathrm{Hom}_{k}(P^{\circ}, M) = P^{k}_{M}.
\]

\pagerange{5}{5}
\section*{Produit tensoriel de catégories $k$-linéaires (page 5)}

\emph{Deux structures sur $P^{k}_{M}$.} Posant $Q = P^{\circ}$, les deux formules
précédentes donnent\footnote{La page étiquette le membre de gauche
$\mathcal{Q}_{M}$ et celui de droite $\mathcal{P}_{M}$. Avec la convention de
la page 9, $P^{k}_{M} = \mathrm{Hom}_{k}(P^{\circ}, M)$, les deux sont
$Q^{k}_{M}$ ; on suit la page 9.}, pour $M$ complète et cocomplète,
\[
\mathrm{Hom}^{!}_{k}\bigl((Q^{k})^{\circ}, M\bigr) \simeq Q^{k}_{M}
= \mathrm{Hom}_{k}(P, M) \simeq \mathrm{Hom}_{k!}(P^{k}, M).
\]
Sur $P^{k}_{M} = \mathrm{Hom}_{k}(P^{\circ}, M)$ opèrent
$(P^{\circ})^{k}$, par $(F, L') \mapsto F *_{k} L'$, et $P^{k}$, par
$(L, F) \mapsto \mathrm{Hom}_{k}(L, F)$, et l'on a
\[
(F *_{k} L')^{\circ} \simeq \mathrm{Hom}_{k}(L', F^{\circ}),
\qquad
\mathrm{Hom}_{k}(L, F)^{\circ} \simeq F^{\circ} *_{k} L ,
\]
le premier dans $M^{\circ}$, où $F^{\circ} : P \to M^{\circ}$.

\textbf{Définition.} Pour $P$, $Q$ $k$-additives, le produit tensoriel
$P \otimes_{k} Q$, muni d'un foncteur $k$-bilinéaire $P \times Q \to P \otimes_{k}
Q$, est défini par la propriété 2-universelle
\[
\mathrm{Hom}_{k}(P \otimes_{k} Q, M) \simeq \mathrm{Bil}_{k}(P, Q; M)
\qquad \text{pour toute } M \ k\text{-additive}.
\]

\emph{Construction explicite.} Les objets sont les familles finies
$\zeta = \bigoplus_{i \in I} a_{i} \otimes_{k} b_{i}$, et
\[
\mathrm{Hom}(\zeta, \zeta') = \prod_{(i,j) \in I \times J}
\mathrm{Hom}_{P}(a_{i}, a'_{j}) \otimes_{k} \mathrm{Hom}_{Q}(b_{i}, b'_{j}),
\]
la composition étant celle des matrices. À retenir :
\begin{enumerate}
\item[a)] si $P$ et $Q$ sont petites, $P \otimes_{k} Q$ l'est ;
\item[b)] $\mathrm{Hom}_{P} (a, a') \otimes_{k} \mathrm{Hom}_{Q}(b, b')
\xrightarrow{\;\sim\;} \mathrm{Hom}_{P \otimes_{k} Q}(a \otimes_{k} b, a'
\otimes_{k} b')$.
\end{enumerate}

\pagerange{7}{7}
\section*{Plongement, et deux applications (page 7)}

Le plongement de Yoneda de $P \otimes_{k} Q$ se lit dans les foncteurs
bilinéaires :
\[
P \otimes_{k} Q \hookrightarrow (P \otimes_{k} Q)^{k} \simeq
\mathrm{Bil}_{k}(P^{\circ}, Q^{\circ}; Ab_{k}),
\qquad
a \otimes_{k} b \longmapsto \bigl((x, y) \mapsto \mathrm{Hom}_{P}(x, a)
\otimes_{k} \mathrm{Hom}_{Q}(y, b)\bigr).
\]
Il est pleinement fidèle par b) et Yoneda. D'où une seconde construction de
$P \otimes_{k} Q$ : la sous-catégorie pleine de
$\mathrm{Bil}_{k}(P^{\circ}, Q^{\circ}; Ab_{k})$ formée des sommes finies de tels
foncteurs --- ce qui prouve l'existence, et la petitesse si $P$ et $Q$ sont
petites.

\textbf{Application 1.} Soit $M \simeq Q^{k}$, et $M^{\vee} \simeq
(Q^{\circ})^{k}$ sa duale. Alors $P^{k}_{M}$ et $(P^{\circ})^{k}_{M^{\vee}}$ sont
encore des catégories de modules :
\[
P^{k}_{M} = \mathrm{Hom}_{k}\bigl(P^{\circ}, \mathrm{Hom}_{k}(Q^{\circ},
Ab_{k})\bigr) \simeq \mathrm{Hom}_{k}(P^{\circ} \otimes_{k} Q^{\circ}, Ab_{k}) =
(P \otimes_{k} Q)^{k},
\]
\[
(P^{\circ})^{k}_{M^{\vee}} \simeq \mathrm{Hom}_{k}(P \otimes_{k} Q, Ab_{k}) =
\bigl((P \otimes_{k} Q)^{\circ}\bigr)^{k},
\]
en utilisant $(P \otimes_{k} Q)^{\circ} = P^{\circ} \otimes_{k} Q^{\circ}$.

\textbf{Application 2.} Pour $M$ $k$-additive cocomplète,
\[
\mathrm{Bil}_{k!!}(P^{k}, Q^{k}; M) \simeq \mathrm{Hom}_{k!}\bigl((P \otimes_{k}
Q)^{k}, M\bigr) \simeq \mathrm{Bil}_{k}(P, Q; M) :
\]
$(P \otimes_{k} Q)^{k}$ est le produit tensoriel de $P^{k}$ et $Q^{k}$ parmi les
catégories $k$-additives cocomplètes et les foncteurs qui commutent aux limites
inductives.\footnote{C'est le produit tensoriel de catégories $k$-linéaires
cocomplètes de la théorie des catégories enrichies (G.~M.~Kelly, 1982), dont
celui de P.~Deligne pour les catégories abéliennes $k$-linéaires de longueur
finie (1990) est un cousin ; les noms ne sont pas sur la page, et les feuillets
ne citent personne. La page pose un point
d'interrogation sous l'énoncé, que la démonstration de la page 9 lève.}

\pagerange{9}{9}
\section*{La curryfication, et la comparaison $(\circledast)$ (page 9)}

\emph{Démonstration de l'application 2.}
\[
\begin{aligned}
\mathrm{Bil}_{k!!}(P^{k}, Q^{k}; M)
 &\simeq \mathrm{Hom}_{k!}\bigl(P^{k}, \mathrm{Hom}_{k!}(Q^{k}, M)\bigr)
  \simeq \mathrm{Hom}_{k!}\bigl(P^{k}, \mathrm{Hom}_{k}(Q, M)\bigr) \\
 &\simeq \mathrm{Hom}_{k}\bigl(P, \mathrm{Hom}_{k}(Q, M)\bigr)
  \simeq \mathrm{Hom}_{k}(P \otimes_{k} Q, M)
  \simeq \mathrm{Hom}_{k!}\bigl((P \otimes_{k} Q)^{k}, M\bigr),
\end{aligned}
\]
en utilisant que $\mathrm{Hom}_{k}(Q, M)$ est cocomplète, les limites
inductives s'y calculant point par point. Dualement,
$\mathrm{Bil}^{!!}_{k}((P^{k})^{\circ}, (Q^{k})^{\circ}; M) \simeq
\mathrm{Hom}^{!}_{k}(((P \otimes_{k} Q)^{k})^{\circ}, M)$ pour $M$ complète.

\textbf{Application 3.} Soit $M$ $k$-additive cocomplète. On a un foncteur
$k$-bilinéaire, d'où un foncteur
\[
(\circledast) \qquad P^{k} \otimes_{k} M \longrightarrow
\mathrm{Hom}_{k}(P^{\circ}, M) = P^{k}_{M},
\qquad
L \otimes_{k} x \longmapsto \bigl(a \mapsto L(a) \otimes_{k} x\bigr) =
\widetilde{x} \circ L,
\]
où $\widetilde{x} : Ab_{k} \to M$, $N \mapsto N \otimes_{k} x$, est le foncteur
cocontinu qui correspond à $x$ par $\mathrm{Hom}_{k!}(Ab_{k}, M) \simeq M$.

\emph{Ce foncteur n'est pas pleinement fidèle en général.} Pour $P = k$ vu comme
catégorie à un objet, $k$ un corps et $M = Ab_{k}$, il s'agit de l'application
\[
\mathrm{Hom}_{k}(L, L') \otimes_{k} \mathrm{Hom}_{k}(x, x') \longrightarrow
\mathrm{Hom}_{k}(L \otimes_{k} x, L' \otimes_{k} x'),
\]
qui n'est pas surjective dès que $L$ et $x$ sont de dimension infinie. Soit
$L = x = L' = x' = V$ de base $(e_{i})_{i \in \mathbf{N}}$, et $D$ la projection
de $V \otimes V$ sur le sous-espace engendré par les $e_{i} \otimes e_{i}$.
Pour un opérateur $T$ de $V \otimes V$, notons $T_{ij} : V \to V$,
$v \mapsto (e_{i}^{*} \otimes \mathrm{id})\,T(e_{j} \otimes v)$. Si
$T = \sum_{r \leq n} A_{r} \otimes B_{r}$, les $T_{ij}$ restent dans l'espace
engendré par $B_{1}, \ldots, B_{n}$ ; or les $D_{ii}$ sont les projecteurs
sur les $e_{i}$, en nombre infini et linéairement indépendants. Il est
pleinement fidèle sur les $L \otimes_{k} x$ avec $L \in \mathrm{Kar}\, P$ lorsque
les $\mathrm{Hom}$ de $P$ sont des $k$-modules projectifs de type fini : pour
$L = a$ représentable, $\mathrm{Hom}(a \otimes_{k} x, L' \otimes_{k} x') =
\mathrm{Hom}_{M}(x, L'(a) \otimes_{k} x')$ par Yoneda, et c'est
$L'(a) \otimes_{k} \mathrm{Hom}_{M}(x, x')$ dès que $L'(a)$ est projectif de
type fini.\footnote{La page affirme « Ce foncteur $(\circledast)$ est pl.\
fidèle. En effet », puis dessine une comparaison par $M \to M^{k}$,
$\mathrm{Hom}_{k}(P^{\circ}, M^{k})$ et $\mathrm{Bil}_{k}(P^{k} \times M^{\circ},
Ab_{k})$, avec la mise en garde « attention, il faut passer à un univers plus
gros ». Elle aboutit à « $P^{k} \otimes_{k} M \dashrightarrow
\mathrm{Bil}_{k}(P^{k} \times M^{\circ}, Ab_{k})$ pl.\ fid.\ ? », en pointillé et
avec un point d'interrogation, et à une remarque dont le premier mot, lu
« Can. », n'est pas assuré. Le contre-exemple et l'énoncé restreint sont les
nôtres.}

\pagerange{11}{11}
\section*{Formules, projectifs, et une question (page 11)}

On identifiera parfois $L \otimes_{k} x$ à son image dans $P^{k}_{M}$. Pour
$M = Q^{k}$, $(\circledast)$ devient
\[
P^{k} \otimes_{k} Q^{k} \longrightarrow (P \otimes_{k} Q)^{k} \simeq
\mathrm{Bil}_{k}(P^{\circ}, Q^{\circ}; Ab_{k}),
\qquad
F \otimes_{k} G \longmapsto \bigl((a, b) \mapsto F(a) \otimes_{k} G(b)\bigr),
\]
avec les mêmes réserves sur la pleine fidélité.\footnote{La page écrit
« pl.\ fid. » sans réserve ; le contre-exemple de la page 9 s'applique, avec
$P = Q = k$.}

\textbf{Formule 1.} Pour $L \in P^{k}$, $x \in M$ et $L' \in (P^{\circ})^{k}$,
\[
(L \otimes_{k} x) *_{k} L' \simeq (L *_{k} L') \otimes_{k} x,
\qquad L *_{k} L' \in Ab_{k},
\]
parce que $\widetilde{x}$ commute aux limites inductives, donc au cofin ; il suffit
de le vérifier pour $L = a$ et $L' = b^{\circ}$ représentables, où les deux
membres valent $\mathrm{Hom}_{P}(b, a) \otimes_{k} x$.

\textbf{Formule 2.} Pour $F \in P^{k}_{M}$,
\[
\mathrm{Hom}_{P^{k}_{M}}(L \otimes_{k} x, F) \simeq
\mathrm{Hom}_{P^{k}}\bigl(L, \mathrm{Hom}_{M}(x, F)\bigr),
\qquad
\mathrm{Hom}_{M}(x, F) : a \longmapsto \mathrm{Hom}_{M}\bigl(x, F(a)\bigr),
\]
et, si $M$ est de plus complète, les deux sont isomorphes à
$\mathrm{Hom}_{M}(x, \mathrm{Hom}_{k}(L, F))$.\footnote{La page écrit la chaîne
des trois termes. Celui du milieu utilise l'objet $\mathrm{Hom}_{k}(L, F)$ de $M$,
qui est une limite projective ; l'isomorphisme des deux termes extrêmes vaut
sans elle :
$\int_{a} \mathrm{Hom}_{M}(L(a) \otimes x, F(a)) = \int_{a}
\mathrm{Hom}_{k}(L(a), \mathrm{Hom}_{M}(x, F(a)))$. Une note en biais dans la
marge donne le cas $L = a$, où les deux derniers membres valent
$\mathrm{Hom}_{M}(x, F(a))$.}

\textbf{Corollaire.} Si $L$ est projectif dans $P^{k}$ et $x$ projectif dans $M$,
$L \otimes_{k} x$ est projectif dans $P^{k}_{M}$ : par la formule 2,
$\mathrm{Hom}(L \otimes_{k} x, -)$ est le composé de $F \mapsto
\mathrm{Hom}_{M}(x, F(-))$, exact puisque $x$ est projectif et que l'exactitude
se lit point par point, et de $\mathrm{Hom}_{P^{k}}(L, -)$, exact.

\textbf{Question.} Si $M$ est cocomplète et a assez de projectifs, $P^{k}_{M}$
a-t-elle assez de projectifs $F$, et ceux-ci annulent-ils les
$\mathrm{Tor}_{i}(L, F)$, $L \in P^{k}$ ?

\pagerange{12}{12}
\section*{Le produit dérivé (page 12)}

La première moitié de la question a une réponse positive, et la seconde une
réponse positive sous une hypothèse de platitude ; la page le dit, dans ses
deux marges.

\textbf{Proposition.} Soit $N$ une catégorie $k$-additive abélienne admettant
les petites sommes.
\begin{enumerate}
\item[a)] Si tout objet de $N$ est quotient d'un projectif, tout $F \in
P^{k}_{N}$ est quotient d'une somme
$\Phi = \bigoplus_{a \in \mathrm{Ob}\, P} a \otimes_{k} x_{a}$, les $x_{a}$
projectifs, et un tel $\Phi$ est projectif.
\item[b)] Si de plus, pour tout projectif $x$ de $N$, le foncteur
$Ab_{k} \to N$, $V \mapsto V \otimes_{k} x$, est exact --- $x$ est
« $k$-plat » ---, alors pour tout projectif $\Phi$ de $P^{k}_{N}$ le foncteur
$L' \mapsto \Phi *_{k} L'$ est exact, et le produit dérivé
$F \circledast^{L}_{k} L'$ se calcule indifféremment par résolutions
projectives de $F$ ou de $L'$.
\end{enumerate}

\emph{Démonstration.} a) Pour chaque $a$, choisissons un épimorphisme
$x_{a} \to F(a)$ ; par Yoneda il correspond à $a \otimes_{k} x_{a} \to F$, et la
somme de ces flèches est un épimorphisme, point par point. $\Phi$ est
projectif par le corollaire de la page 11. b) Pour $\Phi = a \otimes_{k} x$,
la formule 1 donne $\Phi *_{k} L' = L'(a) \otimes_{k} x$, exact en $L'$ par
platitude ; un projectif quelconque est facteur direct d'une somme de tels
objets.\footnote{Les deux démonstrations sont les nôtres. La page énonce
aussi, en tête, que le produit dérivé se définit par résolutions projectives en
$F$ et qu'un quasi-isomorphisme en $L'$ en donne un ; et, dans une liste de
quatre hypothèses sur $N$ dont la deuxième est illisible, elle nomme la
quatrième « $k$-plat ». Elle écrit à deux reprises « $M \mapsto M \otimes_{k} a$ »
là où la phrase parle de $x$ ; on lit $x$. Le signe $\circledast$ est son
astérisque cerclé, qui désigne ici l'opération notée $*_{k}$ à la page 11.
Dans la langue d'aujourd'hui, $F \circledast^{L}_{k} L'$ est le produit
tensoriel dérivé au-dessus de l'« anneau à plusieurs objets » $P$.}

\end{document}
